% ======================================================================
% Chapter 4 -- Part C: build, deploy, evaluate
% ======================================================================

\section{Build: How Do We Make It?}
\label{sec:policy-build}

The build changes the sample in four places: the documents and their metadata, the agent's instruction, the search tool, and a guardrail. The full code is in the companion repository under \texttt{ch04/}; the listings here show the parts that carry design decisions.

\subsection{Documents with metadata}

Filtering needs fields to filter on. The grounded design switches the data connector from unstructured content to structured documents: each policy file is listed in a JSON Lines manifest that gives the file's location and its metadata. Listing~\ref{lst:jsonl} shows one entry.

\begin{lstlisting}[language={}, caption={One entry in the policy manifest (\texttt{policies.jsonl}). The connector's data schema becomes \texttt{document}.}, label={lst:jsonl}]
{"id": "UK-LEAVE-v7",
 "structData": {"title": "Annual Leave (United Kingdom)", "country": "UK",
                "audience": "all", "grade_band": "all", "version": 7,
                "effective_from": "2026-04-01", "owner": "people-policy"},
 "content": {"mimeType": "application/pdf",
             "uri": "gs://af-policy-library/uk/annual-leave-v7.pdf"}}
\end{lstlisting}

The document ID carries the version, so a citation identifies exactly which text an answer relied on. Grade rules are stored as bands (\texttt{all}, \texttt{g7plus}) rather than numbers, so that every filter condition has the same simple form. The manifest is generated by the publishing workflow, which refuses to publish a document whose metadata fails validation---the change that came out of the failure walk-through.

\subsection{The instruction}

\begin{lstlisting}[language={}, caption={The grounded design's instruction (abridged).}, label={lst:policy-instr}]
You answer questions from Alder & Finch employees about company policy.

Rules:
- Answer ONLY from passages returned by search_policies in this conversation.
  Never use general knowledge about employment law or other companies.
- The passages are already limited to policies that apply to this employee.
  Do not ask for their country or grade.
- After every factual sentence, cite its source as [DOC-ID §section],
  for example [UK-LEAVE-v7 §2.1].
- If the passages do not answer the question, or answer it only partly,
  reply exactly: NO_SOURCE
- You may call search_policies up to three times, for example once for a
  general rule and once for an exception. Then answer or reply NO_SOURCE.
- Be brief. Quote numbers exactly as the policy states them.
\end{lstlisting}

Each rule maps to a gap in the sample design or a risk: the first two to \risk{Wrong version}, the citation rule to \kpi{Answer accuracy} and auditability, the \texttt{NO\_SOURCE} rule to \risk{Silent policy}. Note what the instruction does \emph{not} do: it does not enforce the search limit or the permission filter. Those live in code.

\subsection{The permission-aware search tool}

The sample's built-in search tool is configured once, when the agent starts, and is the same for every user. The grounded design needs a filter that depends on who is asking. The solution is a function tool that reads the employee's profile from the session state, builds the filter, and calls the search service (Listing~\ref{lst:search-tool}).

\begin{lstlisting}[language=Python, caption={Profile loading and the permission-aware search tool (\texttt{ch04/app/tools.py}, abridged).}, label={lst:search-tool}]
from google.adk.agents.callback_context import CallbackContext
from google.adk.tools import ToolContext
from .directory import get_profile      # read-only call to the HR directory server
from .search_backend import search      # Agent Search query; returns passages

MAX_SEARCHES = 3

def load_profile(callback_context: CallbackContext) -> None:
    """before_agent_callback: store the signed-in employee's profile in state."""
    if "user:employee" not in callback_context.state:
        email = callback_context.state["user:verified_email"]   # set at sign-in
        callback_context.state["user:employee"] = get_profile(email)

def build_filter(emp: dict) -> str:
    countries = f'"{emp["country"]}", "GLOBAL"'
    audiences = '"all", "managers"' if emp["is_manager"] else '"all"'
    bands = '"all", "g7plus"' if emp["grade"] >= 7 else '"all"'
    return (f"country: ANY({countries}) AND audience: ANY({audiences}) "
            f"AND grade_band: ANY({bands})")

def search_policies(query: str, tool_context: ToolContext) -> dict:
    """Search the Alder & Finch policies that apply to the current employee.

    Args:
        query: A short search query, e.g. "annual leave carry over".

    Returns:
        A dict with 'passages': each has doc_id, section, version,
        effective_from and text. Or a dict with 'error'.
    """
    emp = tool_context.state.get("user:employee")
    if not emp:                                   # fail closed
        return {"error": "No employee profile; search is not allowed."}
    used = tool_context.state.get("temp:searches", 0)
    if used >= MAX_SEARCHES:
        return {"error": "Search limit reached. Answer or reply NO_SOURCE."}
    tool_context.state["temp:searches"] = used + 1

    passages = search(query, filter=build_filter(emp), top_k=5)
    seen = tool_context.state.get("temp:retrieved", [])
    tool_context.state["temp:retrieved"] = seen + [p["doc_id"] for p in passages]
    return {"passages": passages}
\end{lstlisting}

Four details in this listing are design decisions, not coding style:
\begin{itemize}[itemsep=1pt]
  \item The profile comes from a verified identity set at sign-in, never from the conversation. If the model or the user could supply the country, the filter would be only as trustworthy as the prompt.
  \item The tool \textbf{fails closed}: with no profile, it refuses to search rather than searching without a filter.
  \item The search limit is counted in code. The instruction mentions it so that the model plans within it; the code enforces it.
  \item The tool records which documents it returned. The citation guardrail needs that list, and so does the evaluation.
\end{itemize}

\begin{keyidea}
Who is asking comes from a verified sign-in, never from the conversation. If the model or the user could supply the country, the filter would be only as trustworthy as the prompt.
\end{keyidea}

\subsection{The citation guardrail}

Listing~\ref{lst:cite-guard} runs after each model response. If the response is a final answer, it checks that the answer cites at least one source and that every cited document was retrieved in this conversation. Otherwise it replaces the answer with the offer to open an HR ticket.

\begin{lstlisting}[language=Python, caption={The citation guardrail as an after-model callback (\texttt{ch04/app/guards.py}, abridged).}, label={lst:cite-guard}]
import re
from google.adk.models import LlmResponse
from google.genai import types

CITE = re.compile(r"\[([A-Z]{2,6}-[A-Z]+-v\d+) §[\d.]+\]")
OFFER = ("I couldn't find a clear answer in the policies that apply to you. "
         "Would you like me to open an HR ticket with your question?")

def citation_check(callback_context, llm_response: LlmResponse):
    parts = (llm_response.content.parts if llm_response.content else None) or []
    if not parts or any(p.function_call for p in parts):
        return None                                # not a final answer
    text = "".join(p.text or "" for p in parts).strip()
    cited = set(CITE.findall(text))
    retrieved = set(callback_context.state.get("temp:retrieved", []))
    if text == "NO_SOURCE" or not cited or not cited <= retrieved:
        callback_context.state["temp:offer_ticket"] = True
        return LlmResponse(content=types.Content(
            role="model", parts=[types.Part(text=OFFER)]))
    return None                                    # keep the answer
\end{lstlisting}

The agent is then assembled as in the sample, with the new instruction, \texttt{search\_policies} as its only search tool, \texttt{load\_profile} as a before-agent callback and \texttt{citation\_check} as an after-model callback. The ticket tool is a separate step outside the agent, called by the web front end only after the employee clicks to confirm; the agent itself has no write tools at all.

% ----------------------------------------------------------------------
\section{Deploy: How Do We Run It Safely?}
\label{sec:policy-deploy}

This is the book's first deployment, so we take it step by step. The technical focus of the case is getting the basics right once, so that later cases can build on them.

\subsection{Choosing where it runs}

Chapter~\ref{ch:foundations} described three targets. For this pilot the team chose Cloud Run, for two reasons. First, identity: Alder \& Finch already protects internal web applications with Identity-Aware Proxy (IAP), which signs users in with the company's identity provider and passes a verified identity to the service. A Cloud Run service behind IAP fits that pattern without new infrastructure. Second, the HR directory and ticketing MCP servers\index{Model Context Protocol (MCP)!server} will also run on Cloud Run, so the whole solution uses one operational model. Agent Runtime would have been a reasonable alternative, with managed sessions as its main advantage; Chapter~\ref{ch:case-loan} compares the targets in full.

\subsection{Steps}

\begin{enumerate}[itemsep=2pt]
  \item \textbf{Region and data residency.} Employee questions are personal data. With the data protection officer, the team chose a European region for the agent, the search data store and logs, so that German and UK employees' data stays in Europe. Note that the sample defaults to a \texttt{global} location for the data store; that default is a decision to revisit, not to inherit.
  \item \textbf{A dedicated service account} with only the permissions the solution needs: query access to the one search data store, permission to call the models, and permission to invoke the HR directory and ticketing services. It cannot write to the policy bucket or the index.
  \item \textbf{Pinned configuration.} The model name is set to a specific version, not a \texttt{latest} alias. A model change becomes a deliberate configuration change that must pass the evaluation in Section~\ref{sec:policy-evaluate} first.
  \item \textbf{Deploy the agent} with ADK's command, run from the project folder:
\begin{lstlisting}[language=bash, numbers=none]
adk deploy cloud_run --project=$PROJECT --region=$REGION \
    --service_name=policy-assistant ./app
\end{lstlisting}
  \item \textbf{Close public access.} Configure the service to refuse unauthenticated requests and put it behind IAP, limited to the employee directory group. The front end reads the verified e-mail from IAP and creates each session with it in state (\texttt{user:verified\_email}).
  \item \textbf{Event-driven freshness.} Create an Eventarc trigger on the policy bucket for newly written objects. It runs a small Cloud Run job that imports the changed document, adapted from the import script that ships with the sample. Keep the connector's daily sync as a backstop.
  \item \textbf{Cost controls.} Set a maximum number of Cloud Run instances, a billing budget with alerts, and the search-count limit already enforced in the tool.
\end{enumerate}

\begin{watchout}[Watch out: defaults are decisions]
A global data-store location, a \texttt{latest} model alias and a daily refresh are all defaults that someone chose for a different problem. Each one in this list was examined and replaced; none was simply inherited.
\end{watchout}

\begin{designnote}
The sample's README points to the deprecated agent-starter-pack tooling as the old path and recommends ADK's own deploy commands. Expect such shifts during the life of this book. The design decisions in the steps above---residency, least privilege, pinned models, private access, event-driven freshness---do not change when the commands do.
\end{designnote}

% ----------------------------------------------------------------------
\section{Evaluate: How Do We Know It Works?}
\label{sec:policy-evaluate}

\subsection{Measure retrieval separately from answers}

A retrieval-augmented system can fail in two places: retrieval can miss the right passage, or generation can misuse a passage it was given. If you measure only final answers, you cannot tell which happened, and you will fix the wrong thing. The evaluation therefore measures five things, each tied to a step and a KPI:
\begin{description}[leftmargin=1em,itemsep=2pt]
  \item[Recall@$k$] For answerable questions, the share for which a correct document is among the top $k$ retrieved. Measures the index and the filter. Feeds \kpi{Answer accuracy}.
  \item[Mean reciprocal rank (MRR)] The average of $1/\text{rank}$ of the first correct document. A falling MRR with steady recall means right documents are slipping down the list, often a sign of competing near-duplicates.
  \item[Citation precision] The share of citations that refer to retrieved documents. The guardrail should hold this at 1.0; if it does not, the guardrail is broken.
  \item[Abstention] On questions the policies do not answer, the share where the assistant abstained (abstention recall, for \risk{Silent policy}); on answerable questions, the share where it abstained anyway (false abstention, which costs \kpi{Resolved without HR}).
  \item[Leakage] The number of cases in which any document the employee may not see was retrieved. Not answered---\emph{retrieved}. A restricted passage that reaches the model has already left its boundary, whatever the model does with it.
\end{description}

Listing~\ref{lst:rag-metrics} computes these from evaluation records.

\lstinputlisting[language=Python, caption={Retrieval and answer metrics (\texttt{ch04/rag\_metrics.py}).}, label={lst:rag-metrics}, firstline=2, lastline=38]{code/rag_metrics.py}

\begin{thinkabout}
The repository's version includes a four-case demonstration. Look at its output below before reading the paragraph after it. Would you ship this system?
\end{thinkabout}
\begin{lstlisting}[language={}, numbers=none, frame=none, xleftmargin=1.5em]
recall@5            1.00
MRR                 0.75
citation precision  1.00
abstention recall   1.00   false abstention 0.00
cases with leakage  1
\end{lstlisting}
Every answer metric is perfect, and the assistant correctly declined the question about performance calibration. Yet one case leaked: a manager-only document was retrieved for a non-manager. An evaluation that looked only at answers would have reported a flawless system with a confidentiality breach in it.

\begin{keyidea}
Measure leakage at retrieval, not in the answer. A restricted passage that reaches the model has already left its boundary, whatever the model does with it.
\end{keyidea}

\subsection{The evaluation set}

Table~\ref{tab:policy-evalset} describes the 240-question set built for the pilot. Most questions were drawn from real tickets, rewritten to remove personal details, and labeled by HR generalists with the correct documents and answer. The trap categories were written deliberately to probe the risks.

\begin{table}[htbp]
\centering
\small
\caption{Composition of the policy evaluation set.}
\label{tab:policy-evalset}
\begin{tabularx}{\textwidth}{@{}>{\raggedright\arraybackslash}p{3.6cm}r>{\raggedright\arraybackslash}X@{}}
\toprule
\rowcolor{coHead} Category & Questions & What it tests \\
\midrule
Answerable, by country & 120 & Recall, correctness; 40 per country, asked as employees of that country \\
Grade-dependent & 30 & Filtering by grade band; same question asked at two grades \\
Policy silent & 30 & Abstention (\risk{Silent policy}) \\
Permission traps & 30 & Leakage (\risk{Restricted leak}): non-managers asking about manager-only topics \\
Freshness traps & 20 & Questions on policies changed in the last month (\risk{Wrong version}) \\
Multi-hop & 10 & Rule plus exception, needing two searches \\
\midrule
Total & 240 & \\
\bottomrule
\end{tabularx}
\end{table}

\subsection{Thresholds and judging}

The thresholds in Table~\ref{tab:policy-thresholds} are part of the go/no-go\index{go/no-go criteria} criteria for the pilot and are rerun before every change of model, instruction or index configuration. Correctness of the final answer is judged two ways: by a model-based judge using a rubric that compares the answer with HR's reference answer, and by HR reviewers on a random 60 questions. The model judge is used for routine regression runs only after its verdicts agree with the human reviewers' on at least 90\% of the 60.

\begin{table}[htbp]
\centering
\small
\caption{Evaluation thresholds for the policy assistant.}
\label{tab:policy-thresholds}
\begin{tabular}{@{}llr@{}}
\toprule
\rowcolor{coHead} Metric & Step & Threshold \\
\midrule
Recall@5 (answerable) & Index, filter & $\ge$ 0.95 \\
Answer correctness (answerable) & Agent & $\ge$ 0.97 \\
Citation precision & Citation check & $=$ 1.00 \\
Abstention recall (policy silent) & Agent, check & $\ge$ 0.90 \\
False abstention (answerable) & Agent, check & $\le$ 0.08 \\
Leakage (permission traps) & Filter & $=$ 0 \\
Freshness (publish to answerable) & Import on publish & $\le$ 1 hour \\
\bottomrule
\end{tabular}
\end{table}

Running the same set against the sample design is part of the lab. Without metadata, the sample design cannot pass the leakage threshold by construction, and the freshness traps fail whenever a change lands between daily syncs. The informative comparison is on the 120 country questions, where the sample design's errors show how often unfiltered retrieval picks another country's policy---the same error the HR audit found in human answers.

\subsection{Evaluating without watching employees}

The works council's concern (Section~\ref{sec:policy-kpis}) shapes evaluation in production. The weekly accuracy review samples answers with the employee's identity removed, keeps conversation logs for 30 days, and reports only aggregate statistics. Questions that HR reviewers mark as wrong are rewritten without personal details and added to the evaluation set, so that the set grows from real failures without retaining anyone's history.
